\documentclass[12pt,a4paper,fontset=none]{ctexart}
\usepackage[margin=1in]{geometry}
\usepackage{setspace}
\usepackage{booktabs}
\usepackage{longtable}
\usepackage{array}
\usepackage{amsmath,amssymb}
\usepackage{hyperref}
\usepackage{enumitem}

% English text: prefer Times New Roman, fall back to a Times-compatible font.
\IfFontExistsTF{Times New Roman}
  {\setmainfont{Times New Roman}}
  {\setmainfont{TeX Gyre Termes}}

% This Linux environment does not ship the Windows BiauKai font, so we use
% the installed AR PL KaitiM Big5 as the closest Kai-style fallback.
\setCJKmainfont[AutoFakeBold=2,AutoFakeSlant=.2]{AR PL KaitiM Big5}
\setCJKsansfont[AutoFakeBold=2,AutoFakeSlant=.2]{AR PL KaitiM Big5}
\setCJKmonofont{AR PL KaitiM Big5}
\renewcommand{\contentsname}{目錄}

\setstretch{1.2}
\setlength{\parskip}{0.5em}
\setlength{\parindent}{2em}
\newcolumntype{L}[1]{>{\raggedright\arraybackslash}p{#1}}

\title{How Textual Features Influence Social Media Engagement:\\Evidence from Booking.com on X\\[0.5em]
\large Booking.com 社群文字特徵與互動成效之研究提案（雙語版）}
\author{}
\date{\today}

\begin{document}
\maketitle
\tableofcontents
\newpage

\section*{Part I. English Version}
\addcontentsline{toc}{section}{Part I. English Version}

\section{Background and Motivation}
Social media engagement has become an important performance indicator for brands because likes, comments, and shares reflect audience attention, interaction, and the social diffusion of branded content. For a travel platform such as Booking.com, engagement on X can be especially valuable because travel decisions are socially visible, information-intensive, and often shaped by peer endorsement. Mohammad et al. (2024) argue that social commerce interactions can strengthen credibility and influence downstream behavioral outcomes, which makes engagement a meaningful intermediate outcome for hospitality and travel brands. In this context, understanding which textual features encourage stronger engagement can help Booking.com design more effective posts and communicate more persuasively with potential travelers.

This proposal focuses on the textual design of Booking.com's posts rather than on broad brand strategy. The central idea is that the way a post is written can affect how users react to it. Prior literature suggests that content structure, interactivity cues, emotional tone, and discoverability signals can influence user responses on social media (Moran et al., 2019; van der Harst and Angelopoulos, 2024). Building on that logic, this study examines whether text length, question marks, emotional valence, and hashtags are associated with higher engagement for Booking.com posts on X, while controlling for non-textual post elements such as pictures (Li and Xie, 2020), URLs, and at-mentions.

\section{Research Question}
The research question guiding this study is:

\begin{quote}
How do textual features in Booking.com's posts on X influence social media engagement?
\end{quote}

More specifically, the study asks whether longer posts, more interactive wording through question marks, more positive emotional valence, and a greater number of hashtags are associated with higher engagement, after accounting for other post characteristics. This question is theoretically relevant because it links content design decisions to observable consumer response, and it is managerially relevant because it can guide the brand's day-to-day social media writing strategy.

\section{Hypotheses}
The hypotheses are derived from prior studies on social media engagement and content design. First, longer text may provide richer information, context, and persuasive detail, which can help users understand the message and increase their willingness to engage. Gkikas et al. (2022) show that text characteristics, including length, are relevant predictors of engagement in social media contexts. In a travel setting, longer copy may also better communicate destinations, experiences, and value propositions.

Second, question marks are treated as an interactivity cue. De Vries, Gensler, and Leeflang (2012) argue that questions invite audience participation because they implicitly ask users to respond. For a brand such as Booking.com, a question can encourage users to imagine travel choices, express preferences, or reply with opinions, which should raise engagement.

Third, emotional valence captures whether the wording of a post is more positive or negative. Reddy and Prakash (2024) suggest that more positive sentiment can enhance user interaction and loyalty in social media environments. Because tourism communication often relies on aspiration, excitement, and pleasant imagery, positively valenced language is expected to generate stronger reactions.

Fourth, hashtags may improve discoverability, organize topical meaning, and connect a post to broader conversations. Gkikas et al. (2022) demonstrate that text characteristics, including the number of hashtags, positively impact user engagement rates outperforming other post features. For Booking.com, hashtags may make travel posts easier to locate and more visible within trend-based conversations, thereby improving engagement.

Based on this reasoning, the study proposes the following hypotheses:
\begin{enumerate}[label=H\arabic*.]
    \item Text length is positively associated with social media engagement.
    \item Question marks are positively associated with social media engagement.
    \item Emotional valence is positively associated with social media engagement.
    \item The number of hashtags is positively associated with social media engagement.
\end{enumerate}

\section{Data}
The study uses the cleaned dataset in \texttt{data/Final\_Report\_data\_cleaned (2) (1).xlsx}. The dataset contains 766 posts from Booking.com's official X account. Based on the timestamp field, the observed period runs from November 1, 2021 to November 21, 2022. The dataset includes the original post text, engagement counts, media-related indicators, and multiple preprocessed text fields generated for text analysis.

The dependent variable is a combined engagement measure defined as the sum of likes, comments, and shares. This definition is consistent with the group's current research design and is grounded in the idea that engagement reflects multiple forms of observable user interaction. In the present sample, engagement is highly right-skewed: the median combined engagement is 4, the maximum is 2,474, and 104 posts receive zero combined engagement. This pattern is typical for social media data and has implications for model choice.

The dataset is suitable for the proposed analysis because it already includes the main textual and control variables required for the study, including length, question marks, hashtags, emotional valence, pictures, URLs, at-mentions, and emoji counts.

\section{Measurement of Variables}
This study operationalizes all variables using fields that already exist in the cleaned dataset and the preprocessing logic documented in \texttt{py\_export/single/final\_report\_full.py}. The main variables are defined as follows.

\small
\begin{longtable}{@{}L{0.20\textwidth} L{0.18\textwidth} L{0.36\textwidth} L{0.14\textwidth}@{}}
\toprule
Construct & Role & Operational definition & Dataset field \\
\midrule
\endfirsthead
\toprule
Construct & Role & Operational definition & Dataset field \\
\midrule
\endhead
Social media engagement & Dependent variable & Sum of likes, comments, and shares for each post & Derived from \texttt{like}, \texttt{comment}, \texttt{share} \\
Text length & Main independent variable & Number of cleaned English words after removing URLs, mentions, hashtags, and non-letter symbols, then filtering to recognized English words & \texttt{length} \\
Question mark & Main independent variable & Binary dummy tracking if the post contains at least one question mark after removing URLs (0=No, 1=Yes) & \texttt{question} \\
Emotional valence & Main independent variable & Lexicon-based sentiment score centered around neutral tone; higher values indicate more positive wording & \texttt{valence} \\
Hashtag & Main independent variable & Number of hashtags contained in the post & \texttt{hashtag} \\
Picture & Control variable & Number of pictures or visual media items attached to the post & \texttt{picture} \\
URL & Control variable & Number of URLs contained in the post & \texttt{url} \\
At-mention & Control variable & Number of \texttt{@account} mentions contained in the post & \texttt{at\_mention} \\
Emoji & Extension variable & Number of emoji characters contained in the post & \texttt{emoji} \\
\bottomrule
\end{longtable}
\normalsize

The variable definitions are theoretically motivated and empirically grounded in the data preparation workflow. Text length captures how much verbal information a post provides. The \texttt{length} field is based on cleaned text rather than raw character count, which is appropriate because it better reflects the substantive amount of textual content available to readers. The question-mark variable captures a direct interactive cue and is operationalized as a binary dummy variable, since the vast majority of interrogative posts contain only a single question mark.

Emotional valence is measured with the existing \texttt{valence} field, which is generated through lexicon matching on the cleaned tokens and expresses the positivity or negativity of the language used in each post. This proposal treats valence as a continuous sentiment measure rather than splitting sentiment into separate positive and negative variables at this stage. Hashtags are measured as counts because the number of hashtags may reflect stronger discoverability or stronger topic signaling. The control variables capture non-textual or structural post elements that may affect engagement independently of writing style. Emoji is retained as an extension variable because prior literature suggests that emoji can increase engagement (Ko et al., 2022), but this study prioritizes a cleaner four-hypothesis core model in the main analysis.

\section{Proposed Model Specification}
The baseline empirical model examines how the four focal textual features relate to social media engagement while controlling for other post characteristics:

\begin{equation}
\begin{aligned}
\log(1 + \text{Engagement}_i) = {}& \beta_0 + \beta_1\text{TextLength}_i + \beta_2\text{Question}_i + \beta_3\text{Valence}_i \\
& + \beta_4\text{Hashtag}_i + \beta_5\text{Picture}_i + \beta_6\text{URL}_i + \beta_7\text{AtMention}_i + \varepsilon_i
\end{aligned}
\end{equation}

Conceptually, this specification tests whether the proposed textual features explain variation in engagement above and beyond visual or structural elements in the post. Because the dependent variable is a non-negative count-based sum and the sample shows substantial right skew, the dependent variable is modeled as
\[
\log\bigl(1 + \text{Engagement}_i\bigr)
\]
to reduce the influence of extreme observations and improve interpretability. The conceptual definition of engagement, however, remains the raw sum of likes, comments, and shares.

To extend the model, a secondary specification adds emoji as an additional explanatory variable:

\begin{equation}
\begin{aligned}
\log(1 + \text{Engagement}_i) = {}& \beta_0 + \beta_1\text{TextLength}_i + \beta_2\text{Question}_i + \beta_3\text{Valence}_i \\
& + \beta_4\text{Hashtag}_i + \beta_5\text{Picture}_i + \beta_6\text{URL}_i + \beta_7\text{AtMention}_i \\
& + \beta_8\text{Emoji}_i + \varepsilon_i
\end{aligned}
\end{equation}

This second model is not intended to replace the main hypothesis structure. Instead, it functions as an extension analysis that explores whether emoji contributes explanatory value beyond the four core textual features.

\section{Expected Contributions and Managerial Implications}
This study is expected to contribute to the literature on social media content strategy by applying content-feature theory to a travel platform brand. Much of the prior work demonstrates that content design matters, but the effects may vary across industries and brand contexts. By focusing on Booking.com, this study can show how textual features operate in a tourism-related, information-rich, and experience-oriented setting. The proposal therefore contributes both by testing established ideas in a new context and by linking theory more closely to platform-level content decisions.

From a managerial perspective, the results are expected to help Booking.com refine how it writes posts on X. If the proposed relationships are supported, managers may be encouraged to use more informative text, strategically include questions that invite response, maintain a more positive emotional tone, and use hashtags in ways that improve discoverability. The extension analysis on emoji may also indicate whether visual-linguistic cues are worth incorporating more systematically into social media copy. Overall, the project aims to turn textual features into actionable content guidelines rather than treating engagement as a purely descriptive metric.

\section{Limitations and Next Steps}
Several limitations should be acknowledged at the proposal stage. First, the dependent variable combines likes, comments, and shares using equal weights. This is analytically convenient and consistent with the current project design, but the three actions may not reflect identical levels of user involvement. Second, the study focuses on a single brand, which improves contextual consistency but limits generalizability. Third, the text-cleaning process may introduce some language bias because non-English or mixed-language expressions can be simplified or excluded during preprocessing. Fourth, the proposal is theory-driven and measurement-driven; it does not claim empirical confirmation before the statistical analysis is completed.

The next step is to estimate the proposed models, assess whether the hypothesized signs are supported, and compare the baseline model with the emoji extension model. After that, the final report can expand the descriptive statistics, present regression results, discuss robustness checks, and translate the findings into more specific managerial recommendations.

\section{References (Selected)}
\begin{enumerate}
    \item De Vries, L., Gensler, S. and Leeflang, P.S.H. (2012). Popularity of brand posts on brand fan pages: An investigation of the effects of social media marketing. \textit{Journal of Interactive Marketing}, 26(2), 83--91.
    \item Gkikas, D. C., Tzafilkou, K., Theodoridis, P. K., Garmpis, A. and Gkikas, M. C. (2022). How do text characteristics impact user engagement in social media posts? Modeling content readability, length, and hashtags number in Facebook. \textit{International Journal of Information Management Data Insights}, 2, 100067.
    \item Ko, E.E., Kim, D. and Kim, G. (2022). Influence of emojis on user engagement in brand-related user-generated content. \textit{Computers in Human Behavior}.
    \item Li, Y. and Xie, Y. (2020). Is a picture worth a thousand words? An empirical study of image content and social media engagement. \textit{Journal of Marketing Research}, 57(1), 1--19.
    \item Mohammad, A.A., Elshaer, I.A., Azazz, A.M., Kooli, C., Algezawy, M. and Fayyad, S. (2024). The influence of social commerce dynamics on sustainable hotel brand image, customer engagement, and booking intentions. \textit{Sustainability}, 16(14), 6050.
    \item Moran, G., Muzellec, L. and Johnson, D. (2019). Message content features and social media engagement: evidence from the media industry. \textit{Journal of Product and Brand Management}, 29(5), 533--545.
    \item Reddy, M.A.S. and Prakash, C. (2024). Sentiment Analysis of Social Media Networking Sites: A Comparative Study on User Engagement and Public Opinion.
    \item van der Harst, J.P. and Angelopoulos, S. (2024). Less is more: Engagement with the content of social media influencers. \textit{Journal of Business Research}, 181, 114746.
\end{enumerate}

\newpage
\section*{Part II. 繁體中文版本}
\addcontentsline{toc}{section}{Part II. 繁體中文版本}

\section{研究背景與動機}
社群媒體互動（engagement）已成為品牌重要的績效指標，因為按讚、留言與分享能反映受眾注意力、互動程度，以及品牌內容在社群中的擴散效果。對 Booking.com 這類旅遊平台而言，X 上的互動更具價值，因為旅遊決策高度社會化、資訊需求高，且常受到同儕意見影響。Mohammad 等人（2024）指出，社群商務互動能提升品牌可信度，並影響後續行為結果，因此對旅宿與旅遊品牌而言，互動可視為具管理意義的中介成果。

本研究聚焦於 Booking.com 貼文的「文字設計」，而非整體品牌策略。核心概念是：貼文怎麼寫，會影響使用者如何反應。既有文獻顯示，內容結構、互動線索、情緒語調與可發現性訊號，都可能影響社群回應（Moran 等人，2019；van der Harst 與 Angelopoulos，2024）。據此，本研究檢驗文字長度、問號、情緒正向程度與 hashtag 是否與更高互動有關，並控制圖片（Li 與 Xie，2020）、URL 與 at-mention 等非文字因素。

\section{研究問題}
本研究的核心問題如下：

\begin{quote}
Booking.com 在 X 上的貼文文字特徵，如何影響社群互動表現？
\end{quote}

更具體地說，本研究探討：在控制其他貼文特徵後，較長文字、較多問號、較正向情緒語氣與較多 hashtag，是否與更高互動顯著相關。此問題在理論上連結了內容設計決策與可觀察的消費者反應；在管理上則可直接指引品牌日常社群文案策略。

\section{研究假說}
本研究假說建構自社群互動與內容設計相關文獻。首先，較長文字可能提供更豐富資訊、脈絡與說服細節，協助受眾理解訊息並提高互動意願。Gkikas 等人（2022）指出，文字特徵（含長度）是社群互動的重要預測因子。在旅遊情境中，較長文案也更有機會完整傳達目的地體驗與價值主張。

其次，問號可視為互動提示。De Vries、Gensler 與 Leeflang（2012）認為，提問會隱含邀請受眾回應，因而促進參與。對 Booking.com 而言，提問可鼓勵使用者想像旅遊選擇、表達偏好或留言回覆，進而提升互動。

第三，情緒正向程度（valence）反映貼文語句偏正面或偏負面。Reddy 與 Prakash（2024）指出，較正向情緒有助於提升社群互動與忠誠。在旅遊溝通中，由於訊息常依賴嚮往、興奮與愉悅想像，正向語氣預期能帶來更強反應。

第四，hashtag 可能提升可發現性、組織主題意義，並將貼文連結至更廣泛對話。Gkikas 等人（2022）實證指出，貼文特徵中的 hashtag 數量對提升社群互動率有顯著的正向貢獻。對 Booking.com 而言，hashtag 有助於旅遊內容被搜尋與被看見，從而提高互動。

基於上述推論，提出以下假說：
\begin{enumerate}[label=H\arabic*.]
    \item 文字長度與社群互動呈正向關係。
    \item 問號使用與社群互動呈正向關係。
    \item 情緒正向程度與社群互動呈正向關係。
    \item hashtag 數量與社群互動呈正向關係。
\end{enumerate}

\section{資料說明}
本研究使用清理後資料集 \texttt{data/Final\_Report\_data\_cleaned (2) (1).xlsx}，共包含 Booking.com 官方 X 帳號 766 筆貼文。依時間戳記欄位，觀察期間為 2021 年 11 月 1 日至 2022 年 11 月 21 日。資料含原始貼文文字、互動數據、媒體指標與多項文字前處理欄位。

依變數為「綜合互動量」，定義為 likes、comments、shares 三者加總。此定義符合目前小組研究設計，並反映互動是多種可觀察行為的集合。樣本中互動分布明顯右偏：中位數為 4、最大值為 2,474，且有 104 筆貼文的綜合互動為 0。此型態符合社群資料常見特徵，也影響模型選擇。

資料已包含本研究主要自變數與控制變數，包括 length、question、hashtag、valence、picture、url、at-mention、emoji，適合直接進行後續模型估計。

\section{變數衡量}
本研究所有變數皆以現有清理資料欄位操作化，並對應 \texttt{py\_export/single/final\_report\_full.py} 的前處理邏輯。主要變數定義如下。

\small
\begin{longtable}{@{}L{0.18\textwidth} L{0.16\textwidth} L{0.40\textwidth} L{0.14\textwidth}@{}}
\toprule
構念 & 角色 & 操作型定義 & 資料欄位 \\
\midrule
\endfirsthead
\toprule
構念 & 角色 & 操作型定義 & 資料欄位 \\
\midrule
\endhead
社群互動量 & 依變數 & 每篇貼文 likes、comments、shares 之加總 & 由 \texttt{like}、\texttt{comment}、\texttt{share} 推導 \\
文字長度 & 主要自變數 & 移除 URL、mentions、hashtags 與非字母符號後，再過濾為可辨識英文詞彙之詞數 & \texttt{length} \\
問號 & 主要自變數 & 虛擬變數（若貼文移除 URL 後含有一或多個問號則為 1，否則為 0） & \texttt{question} \\
情緒正向程度 & 主要自變數 & 詞典法情緒分數，以中性為中心；值越高代表語氣越正向 & \texttt{valence} \\
hashtag 數量 & 主要自變數 & 貼文中 hashtag 出現次數 & \texttt{hashtag} \\
圖片 & 控制變數 & 貼文附帶圖片或視覺媒體數量 & \texttt{picture} \\
URL & 控制變數 & 貼文中 URL 數量 & \texttt{url} \\
at-mention & 控制變數 & 貼文中 \texttt{@account} 提及次數 & \texttt{at\_mention} \\
emoji & 延伸變數 & 貼文中 emoji 字元數量 & \texttt{emoji} \\
\bottomrule
\end{longtable}
\normalsize

上述變數設定兼具理論與實務可行性。\texttt{length} 以清理後詞數而非原始字元數衡量，較能反映讀者可接收的實質文字資訊量。考量樣本中絕大多數含疑問語氣之貼文僅有單一問號，具有高度偏態，因此問號特徵以虛擬變數（dummy variable）處理，將重點聚焦於互動意圖的有無。\texttt{valence} 以連續分數呈現語言正負向程度，暫不拆分為正面與負面兩個獨立構面。hashtag 以計數表示主題標記與可發現性強度。控制變數用以捕捉文字以外的結構因素。基於文獻指出 emoji 有助於提升互動（Ko 等人，2022），本研究將 emoji 保留於延伸模型中，作為額外解釋力檢驗。

\section{模型設定}
基準模型檢驗四個核心文字特徵對互動的影響，並控制其他貼文特徵：

\begin{equation}
\begin{aligned}
\log(1 + \text{Engagement}_i) = {}& \beta_0 + \beta_1\text{TextLength}_i + \beta_2\text{Question}_i + \beta_3\text{Valence}_i \\
& + \beta_4\text{Hashtag}_i + \beta_5\text{Picture}_i + \beta_6\text{URL}_i + \beta_7\text{AtMention}_i + \varepsilon_i
\end{aligned}
\end{equation}

此模型用以檢驗：文字特徵是否能在視覺與結構因素之外，額外解釋互動差異。由於依變數為非負計數加總且右偏明顯，本研究以
\[
\log\bigl(1 + \text{Engagement}_i\bigr)
\]
作為實證估計之依變數，以降低極端值影響並提升解讀性；互動的概念定義仍維持 likes、comments、shares 的原始加總。

延伸模型加入 emoji：

\begin{equation}
\begin{aligned}
\log(1 + \text{Engagement}_i) = {}& \beta_0 + \beta_1\text{TextLength}_i + \beta_2\text{Question}_i + \beta_3\text{Valence}_i \\
& + \beta_4\text{Hashtag}_i + \beta_5\text{Picture}_i + \beta_6\text{URL}_i + \beta_7\text{AtMention}_i \\
& + \beta_8\text{Emoji}_i + \varepsilon_i
\end{aligned}
\end{equation}

此延伸模型目的不在取代主假說架構，而是評估 emoji 是否能在四項核心文字特徵之外提供額外解釋力。

\section{預期貢獻與管理意涵}
本研究預期在社群內容策略文獻中提供兩項貢獻。第一，將內容特徵理論應用於旅遊平台品牌情境，檢驗既有理論在不同產業脈絡下是否仍成立。第二，將理論更直接連結到平台層級的內容決策，提升研究對實務的可用性。

就管理意涵而言，若假說獲得支持，Booking.com 可更有策略地調整 X 貼文寫作方式，例如提高資訊量、適度加入提問、維持較正向語氣，以及以提升可發現性為目標配置 hashtag。emoji 的延伸結果亦可協助判斷是否應更系統化納入文字與視覺混合的文案設計。整體而言，本研究目標是把「文字特徵」轉化為可執行的內容準則，而非僅將互動視為描述性指標。

\section{限制與後續規劃}
本研究在提案階段有幾點限制。第一，依變數將 likes、comments、shares 以等權重加總，雖具分析便利性，但三者可能代表不同程度的參與。第二，研究僅聚焦單一品牌，有助於提升情境一致性，但限制外部推論。第三，文字清理流程可能對非英文或混合語言表達產生偏誤。第四，本提案屬理論與測量設計，尚不代表實證結果已被驗證。

下一步將估計上述模型，檢驗假說方向是否成立，並比較基準模型與 emoji 延伸模型表現。其後可於最終報告補充描述統計、迴歸結果、穩健性檢驗，並轉化為更具體的管理建議。

\section{參考文獻（節錄）}
\begin{enumerate}
    \item De Vries, L., Gensler, S. and Leeflang, P.S.H. (2012). Popularity of brand posts on brand fan pages: An investigation of the effects of social media marketing. \textit{Journal of Interactive Marketing}, 26(2), 83--91.
    \item Gkikas, D. C., Tzafilkou, K., Theodoridis, P. K., Garmpis, A. and Gkikas, M. C. (2022). How do text characteristics impact user engagement in social media posts? Modeling content readability, length, and hashtags number in Facebook. \textit{International Journal of Information Management Data Insights}, 2, 100067.
    \item Ko, E.E., Kim, D. and Kim, G. (2022). Influence of emojis on user engagement in brand-related user-generated content. \textit{Computers in Human Behavior}.
    \item Li, Y. and Xie, Y. (2020). Is a picture worth a thousand words? An empirical study of image content and social media engagement. \textit{Journal of Marketing Research}, 57(1), 1--19.
    \item Mohammad, A.A., Elshaer, I.A., Azazz, A.M., Kooli, C., Algezawy, M. and Fayyad, S. (2024). The influence of social commerce dynamics on sustainable hotel brand image, customer engagement, and booking intentions. \textit{Sustainability}, 16(14), 6050.
    \item Moran, G., Muzellec, L. and Johnson, D. (2019). Message content features and social media engagement: evidence from the media industry. \textit{Journal of Product and Brand Management}, 29(5), 533--545.
    \item Reddy, M.A.S. and Prakash, C. (2024). Sentiment Analysis of Social Media Networking Sites: A Comparative Study on User Engagement and Public Opinion.
    \item van der Harst, J.P. and Angelopoulos, S. (2024). Less is more: Engagement with the content of social media influencers. \textit{Journal of Business Research}, 181, 114746.
\end{enumerate}

\end{document}
